%% Diagramme de définition de blocs (bdd) de la navette maritime électro-solaire.
%% Arborescence système -> sous-système -> blocs ; losange plein = composant
%% obligatoire, losange vide = composant optionnel ; le nombre porté sur le lien
%% est le nombre d'éléments identiques.
%% Mise en page « téléphone » : les quatre blocs sont rangés en 2 x 2 ; les deux
%% blocs du bas sont reliés par les côtés (liens qui contournent la grille).
\documentclass[tikz,border=4pt]{standalone}
\usepackage{liaisons}
\begin{document}
\begin{tikzpicture}[x=1cm,y=1cm,
  cadre/.style={draw=black!70, line width=0.9pt},
  lien/.style={line width=0.8pt, black!75},
  boite/.style={draw=solideE, line width=0.9pt, fill=solideE!7, inner sep=2.5pt,
                font=\scriptsize,
                execute at begin node={\hyphenpenalty=10000\exhyphenpenalty=10000}},
  tete/.style={boite, align=center},
  vals/.style={boite, align=left, fill=white},
  mult/.style={font=\scriptsize, inner sep=1.5pt}]

%% ---- macro « bloc » : \bloc{nom TikZ}{x}{y}{largeur}{stéréotype}{nom}{values} ----
\newcommand\bloc[7]{%
  \node[tete, text width=#4, anchor=north] (#1h) at (#2,#3) {\textit{#5}\\[1pt] \textbf{#6}};
  \node[vals, text width=#4, anchor=north] (#1v) at (#1h.south) {#7};}
%% ---- losanges de composition : plein (obligatoire) / vide (optionnel) ----
\newcommand\losangeplein[2]{%
  \fill[black] (#1,#2) -- (#1+0.12,#2-0.16) -- (#1,#2-0.32) -- (#1-0.12,#2-0.16) -- cycle;}
\newcommand\losangevide[2]{%
  \draw[lien, fill=white] (#1,#2) -- (#1+0.12,#2-0.16) -- (#1,#2-0.32) -- (#1-0.12,#2-0.16) -- cycle;}

%% ---- repères de la grille -------------------------------------------
\def\cgauche{2.00}  % centre de la colonne de gauche
\def\cdroite{5.60}  % centre de la colonne de droite
\def\rhaut{3.20}    % haut des blocs de la rangée du haut
\def\rbas{0.95}     % haut des blocs de la rangée du bas
\def\lb{3cm}        % largeur de texte d'un bloc

%% ================== cadre à onglet ====================================
\draw[cadre] (0,-0.55) rectangle (7.6,6.85);
\draw[cadre] (0,6.40) -- (4.5,6.40) -- (4.75,6.63) -- (4.75,6.85);
\node[font=\scriptsize, anchor=west, inner sep=3pt] at (0.1,6.62)
     {\texttt{bdd} [Paquet] Blocks [BDD Ferry Boat]};

%% ================== niveau 1 : le système =============================
\bloc{sys}{3.80}{6.20}{3.3cm}{«system»}{Ferry Boat}
     {Longueur = 13~m ; Largeur = 4,70~m ; Vitesse = 7~nœuds}

%% ================== niveau 2 : le sous-système ========================
\node[tete, text width=2.9cm, anchor=north] (sub) at (3.80,4.30)
     {\textit{«subsystem»}\\[1pt] \textbf{Alimentation électrique}};
\losangeplein{3.80}{4.82}
\draw[lien] (3.80,4.50) -- (3.80,4.30);

%% ================== niveau 3 : les quatre blocs (2 x 2) ===============
\bloc{pv1}{\cgauche}{\rhaut}{\lb}{«block»}{Panneaux photovoltaïques Propulsion}
     {Modèle = CENIT~220\\ Puissance crête = 220~Wc\\ Rendement = 13,8~\%}
\bloc{pv2}{\cdroite}{\rhaut}{\lb}{«block»}{Panneaux photovoltaïques Service}
     {Modèle = CENIT~150\\ Puissance crête = 150~Wc\\ Rendement = 11,8~\%}
\bloc{bat}{\cgauche}{\rbas}{\lb}{«block»}{Parc 1 Batteries Propulsion}
     {Technologie = Ni-Cd\\ Tension = 384~V}
\bloc{ven}{\cdroite}{\rbas}{\lb}{«block»}{Ventilateurs batteries}
     {Modèle = JABSCO DC150CFM}

%% ================== liens de composition ==============================
%% rail horizontal sous le sous-système, porteur des quatre losanges
\draw[lien] (2.90,3.62) -- (4.70,3.62)  (3.80,3.62) -- (3.80,3.72);
\losangeplein{2.90}{3.62} \losangeplein{3.50}{3.62}
\losangeplein{4.10}{3.62} \losangevide{4.70}{3.62}
%% rangée du haut : descente directe sur le haut des blocs
\draw[lien] (3.50,3.30) -- (3.50,3.40) -- (\cgauche,3.40) -- (\cgauche,\rhaut);
\draw[lien] (4.10,3.30) -- (4.10,3.40) -- (\cdroite,3.40) -- (\cdroite,\rhaut);
%% rangée du bas : les liens contournent la grille par les côtés
\draw[lien] (2.90,3.30) -- (2.90,3.52) -- (0.25,3.52) |- (bath.west);
\draw[lien] (4.70,3.30) -- (4.70,3.52) -- (7.35,3.52) |- (venh.east);
\node[mult, anchor=north west] at (2.07,3.18) {16};
\node[mult, anchor=north west] at (5.67,3.18) {8};

%% ================== légende ===========================================
\losangeplein{0.35}{-0.75}
\node[font=\scriptsize, anchor=west] at (0.5,-0.91) {obligatoire};
\losangevide{2.35}{-0.75}
\node[font=\scriptsize, anchor=west] at (2.5,-0.91) {optionnel};
\node[font=\scriptsize, anchor=north west] at (0.25,-1.15)
     {--- le nombre porté sur le lien est le nombre d'éléments identiques};
\end{tikzpicture}
\end{document}
